\chapter{Introduction}
\label{chap:introduction}

Particle physics is dedicated to unraveling the most fundamental constituents of the universe and the forces that govern their interactions. It explores the known elementary particles, such as quarks, leptons, and bosons, which constitute matter and mediate the fundamental forces, and it organizes them within the Standard Model (SM), the central theoretical framework of the field. The SM has been remarkably successful in describing the behavior of these particles, yet the discovery that neutrinos have mass and undergo flavor-changing oscillation showed that the model, in its original massless-neutrino formulation, is necessarily incomplete. This has made the neutrino sector one of the most active windows through which physics beyond the Standard Model (BSM) may first reveal itself.

Neutrinos were theoretically proposed by Pauli in 1930 to preserve energy and angular momentum conservation in beta decay, and they were experimentally discovered in 1956 by Cowan and Reines, after decades in which the technology required to detect their exceedingly weak interaction with matter simply did not exist. Since then, neutrinos have been studied in ever greater detail, from the discovery of their oscillations to the ongoing effort to pin down their absolute mass scale, their Dirac or Majorana nature, and their possible electromagnetic properties. Among these properties, the neutrino charge radius (NCR) occupies a particular place: unlike a magnetic moment or an electric dipole moment, it is not an independent addition to the SM but a genuine radiative prediction of electroweak theory itself, arising unavoidably at one loop even in the absence of any new physics. This makes the NCR simultaneously a precision test of the SM and a potential probe of BSM scenarios that would shift its predicted value away from the SM benchmark.

This work focuses on Coherent Elastic Neutrino-Nucleus Scattering (CE$\nu$NS), a process whose cross section is precisely predicted by the SM and is directly sensitive to the electroweak couplings that the NCR modifies. This sensitivity, however, comes with the fact that CE$\nu$NS events correspond to extremely low-energy nuclear recoils, which are experimentally challenging to observe and which must be related, through a detector-dependent quenching factor, to the ionization energy that is actually measured. Furthermore, the shift induced by the NCR on the effective weak coupling is, at leading order, degenerate with a shift in the weak mixing angle $\sin^2\theta_W$ itself, so that any realistic sensitivity study must treat these two SM quantities together rather than in isolation. In this thesis we study the projected sensitivity of a reactor-based CE$\nu$NS experiment using a liquid argon (LAr) target to both $\sin^2\theta_W$ and the electron neutrino charge radius $\langle r_{\nu_e}^2\rangle$, using as a concrete benchmark the SBC-ININ-LAr proposal exposed to the antineutrino flux of the TRIGA Mark III reactor. We show how the interplay between detector mass, exposure time, and the realistic treatment of the quenching response determines how well these two parameters can be constrained, and to what extent their mutual degeneracy can be resolved.

The thesis introduces, in Chapter 2, the electroweak theory of the SM, beginning with a brief history of the neutrino and its role as the catalyst for the proposal and eventual discovery of the weak nuclear force, followed by the unification of that force with electromagnetism into the electroweak theory that underlies all subsequent calculations in this work. In Chapter 3, we turn to radiative corrections within this framework, showing how the inclusion of one-loop diagrams to the effective neutrino-fermion interaction predicts the electromagnetic properties of the neutrino, and in particular the neutrino charge radius, and how this prediction modifies the SM parameters, most notably $\sin^2\theta_W$, that are probed experimentally. Chapter 4 introduces CE$\nu$NS itself as the neutral-current interaction through which these SM parameters are experimentally accessed, presenting its cross section and combining it with a model of the reactor antineutrino flux to compute the expected number of events for a broadly-described experimental design, a quantity that later chapters use to test the influence of different SM parameter values. Chapter 5 presents the projected sensitivity of the SBC-ININ-LAr experiment to $\sin^2\theta_W$ and $\langle r_{\nu_e}^2\rangle$, for a range of detector masses, exposure times, and quenching factor treatments, together with the corresponding $\Delta\chi^2$ confidence intervals. Finally, we give our conclusions in Chapter 6.